\appendix
\renewcommand{\appendixname}{Lampiran}
\chapter{Tabel Notasi}
\label{app:notasi}

Notasi berikut berlaku di seluruh buku. Bila sebuah bab memakai simbol yang tidak
tercantum di sini, simbol itu diperkenalkan di tempat pemakaiannya. Angka di ujung
setiap butir menunjuk bab yang memakainya.

\section*{Himpunan, logika, dan bilangan (Bab 1--5, 22)}
\begin{description}[style=nextline,itemsep=0.25em]
  \item[$\mathbb{N}$, $\mathbb{Z}$, $\mathbb{Q}$, $\mathbb{R}$, $\mathbb{C}$]
        bilangan asli, bulat, rasional, real, dan kompleks.
  \item[$\forall$, $\exists$, $\neg$] untuk semua, ada, dan bukan.
  \item[$\Rightarrow$, $\Leftrightarrow$] berakibat, dan jika dan hanya jika.
  \item[$\in$, $\notin$, $\subseteq$, $\cup$, $\cap$, $\setminus$] keanggotaan,
        bukan anggota, himpunan bagian, gabungan, irisan, selisih (Bab 2).
  \item[$\mathcal{P}(S)$, $|S|$, $S\times T$] himpunan kuasa, kardinalitas,
        hasilkali kartesius (Bab 2).
  \item[$\sup S$, $\inf S$] batas atas terkecil dan batas bawah terbesar (Bab 3).
  \item[$a_n\to a$] barisan $(a_n)$ konvergen ke $a$ (Bab 4).
  \item[$\sum_{n=1}^{\infty} a_n$] deret dan jumlahnya (Bab 5).
  \item[$\mathrm{fl}(x)$, $\mathrm{ulp}$] hasil pembulatan bilangan real dan satuan
        terkecil yang dapat diwakili (Bab 22).
\end{description}

\section*{Aljabar linear (Bab 6--10)}
\begin{description}[style=nextline,itemsep=0.25em]
  \item[$\mathbf{x}$, $\mathbf{x}^{\top}$, $\|\mathbf{x}\|$] vektor kolom,
        transpos, dan norma Euklides.
  \item[$\operatorname{span} S$, $\dim V$] rentang himpunan vektor dan dimensi
        ruang vektor.
  \item[$\mathcal{L}(V,W)$] ruang semua transformasi linear dari $V$ ke $W$ (Bab 7).
  \item[$\ker T$, $\operatorname{im} T$] inti dan jangkauan transformasi (Bab 7).
  \item[$\det A$, $\operatorname{tr} A$, $A^{-1}$] determinan, jejak, dan invers
        matriks (Bab 8).
  \item[$\lambda$, $\mathbf{v}$] nilai eigen dan vektor eigen (Bab 9).
  \item[$\sigma_i$, $A=U\Sigma V^{\top}$] nilai singular dan dekomposisi nilai
        singular (Bab 10).
  \item[$A^{+}$, $\kappa(A)$] pseudoinvers dan bilangan kondisi (Bab 10, 20).
\end{description}

\section*{Analisis dan kalkulus (Bab 11--16)}
\begin{description}[style=nextline,itemsep=0.25em]
  \item[$\varepsilon$, $\delta$] bilangan positif kecil dalam definisi limit,
        kekontinuan, dan kekontinuan seragam.
  \item[$f'(x)$, $\nabla f$] turunan biasa dan gradien.
  \item[$J_f$, $H_f$] matriks Jacobian dan Hessian (Bab 12).
  \item[$\int_a^b f$] integral Riemann pada $[a,b]$ (Bab 13).
  \item[$\sum_{\mathcal{P}} M_i \Delta x_i$] jumlah atas untuk partisi
        $\mathcal{P}$ (Bab 13).
  \item[$f_n \rightrightarrows f$] barisan fungsi konvergen seragam (Bab 14).
  \item[$\dfrac{\partial(x,y)}{\partial(u,v)}$] determinan Jacobian perubahan
        variabel (Bab 15).
  \item[$x'=f(t,x)$, $\dot{x}=F(x)$] persamaan diferensial biasa dan sistem
        otonom, dengan $'$ dan $\dot{\ }$ sama-sama menyatakan turunan terhadap
        waktu (Bab 16).
  \item[$\varphi_t(x_0)$] aliran: keadaan pada waktu $t$ bila mulai dari $x_0$
        (Bab 16).
  \item[$x^{*}$, $J_F(x^{*})$] titik setimbang dan matriks Jacobian di titik itu;
        tanda bagian real nilai eigennya menentukan kestabilan (Bab 16).
\end{description}

\section*{Peluang dan statistika (Bab 17--21, 26)}
\begin{description}[style=nextline,itemsep=0.25em]
  \item[$P(A)$, $P(A\mid B)$] peluang kejadian $A$ dan peluang bersyarat.
  \item[{$\mathbb{E}[X]$, $\operatorname{Var}X$}] harapan dan variansi.
  \item[$\operatorname{Cov}(X,Y)$] kovarians dua variabel acak (Bab 18).
  \item[$\bar{X}_n$, $\mathrm{se}$] rata-rata sampel dan galat bakunya.
  \item[$\mathrm{CI}_{1-\alpha}$] interval kepercayaan dengan tingkat
        $1-\alpha$ (Bab 20).
  \item[$\mathrm{TP}$, $\mathrm{FP}$, $\mathrm{FN}$, $\mathrm{TN}$] empat sel
        matriks konfusi (Bab 27).
  \item[$\hat{\beta}$, $\hat{y}=X\hat{\beta}$] koefisien regresi kuadrat
        terkecil dan nilai yang diprediksi (Bab 21).
  \item[$e_i=y_i-\hat{y}_i$, $R^2$] residu dan proporsi variansi yang dijelaskan
        model (Bab 21).
  \item[$\mathrm{se}(\hat{\beta}_j)$, $\mathrm{VIF}_j$] galat baku koefisien dan
        faktor pembengkakan variansi (Bab 21).
\end{description}

\section*{Optimasi, pembelajaran mesin, dan verifikasi (Bab 23--30)}
\begin{description}[style=nextline,itemsep=0.25em]
  \item[$\min f_0(x)$ dengan $f_i(x)\le 0$] bentuk baku masalah optimasi (Bab 23).
  \item[$L(x,\lambda,\nu)$] Lagrangian; $\lambda,\nu$ pengali (Bab 23).
  \item[$c^{\top}x$, $Ax\le b$, $x\in\mathbb{Z}^n$] bentuk LP dan program integer
        (Bab 24).
  \item[$\sigma(\cdot)$, $\mathrm{softmax}$] fungsi aktivasi dan normalisasi
        eksponensial (Bab 25).
  \item[$\nabla_{\mathbf{x}}\ell$] gradien fungsi kerugian terhadap masukan;
        bahan bakar serangan (Bab 25, 29).
  \item[$L_2$, $L_1$, $\lambda_{\text{reg}}$] jenis regularisasi dan kekuatannya
        (Bab 26).
  \item[$\|\delta\|_{\infty}\le\varepsilon$] batas gangguan pada serangan
        adversarial (Bab 29, 30).
  \item[$\vdash$, \texttt{by}, \texttt{sorry}] keterbuktian, pembuka blok taktik,
        dan bukti yang ditangguhkan di Lean~4 (Bab 28).
\end{description}
